\documentclass[12pt]{article}
\usepackage[utf8]{inputenc}
\usepackage[T1]{fontenc}
\usepackage{hyperref}
\title{Ayoquesco Zapotec Language Resources for K-12 Curriculum Development}
\author{Compiled by Claude Sonnet 5.5 (claude-sonnet-5-5) for Andr\'as Kornai}
\date{}
\begin{document}
\maketitle

\section{Introduction}
Ayoquesco Zapotec (ISO 639-3 zaf, Glottocode ayoq1235), also called Western Ejutla Zapotec, is spoken in Santa Mar\'ia Ayoquezco and nearby in Oaxaca, Mexico. About 880 speakers are reported, and the language is endangered. Documentation is thin.

\section{Primary Educational Resources}
No primers, textbooks or school curricula were found in the search results.

\section{Linguistic and Reference Materials}
\begin{itemize}
\item Wikipedia article, which cites Robert MacLaury's 1970 master's thesis ``Ayoquesco Zapotec: Ethnography, Phonology, and Lexicon'' as the main early study: \url{https://en.wikipedia.org/wiki/Ayoquezco_Zapotec}
\item Wikitongues language page: \url{https://wikitongues.org/languages/zaf/}
\item Endangered Languages Project entry: \url{https://www.endangeredlanguages.com/node/40413}
\item American Philosophical Society guide to indigenous materials, entry for Ayoquesco Zapotec: \url{https://indigenousguide.amphilsoc.org/entry/39743}
\item Verbix endangered languages page for Ayoquezco: \url{https://docs.verbix.com/EndangeredLanguages/Ayoquezco}
\end{itemize}

\section{Religious and Community Resources}
\begin{itemize}
\item Global Recordings Network audio Bible lessons in 40 sections: \url{https://www.globalrecordings.net/en/language/zaf}
\item Joshua Project profile of the people group: \url{https://joshuaproject.net/index.php/people_groups/15364}
\end{itemize}

\section{Digital and Audio Resources}
The Global Recordings audio is the only audio found.

\section{Curriculum Development Context}
The language is used as a first language by older adults only and is not known to be taught in schools. Any curriculum would be a revitalization effort. It would have to draw on the MacLaury lexicon and on work with elderly speakers.
\end{document}